\documentclass[10pt,a4paper,titlepage,oneside]{report}

% Schriftart
\usepackage[utf8]{inputenc}
\usepackage[german]{babel}
%\usepackage[T1]{fontenc}
%\usepackage[scaled]{uarial} %WINDOWS only!!!
%\renewcommand*\familydefault{\sfdefault}
%\sloppy

\usepackage{graphicx}
\usepackage{epstopdf}


% Deutsches Datumsformat
\usepackage{datetime}
\newdateformat{mydate}{\twodigit{\THEDAY}.\twodigit{\THEMONTH}.\THEYEAR}
\mydate

% Meta Information festlegen
\usepackage[
pdftitle={Ausarbeitung für das Modul Compilerbau 2},
pdfsubject={Tree-Height Balancing},
pdfauthor={Sören Gutzeit, Daniel Vogel, Artjom Siebert, Patrick Koteras},
pdfkeywords={Compilerbau, Tree, Height, Balancing}
]{hyperref}

% Acronyme für Glossar
\usepackage[printonlyused]{acronym}

% Listings
\usepackage{listings}

% biblatex einfügen
\usepackage[babel]{csquotes}
\usepackage[style=alphabetic,hyperref=true,backend=bibtex8]{biblatex}
\bibliography{literatur}


% Eigene Befehle
\newcommand{\acfoot}[1]{\acs{#1}\footnote{\acl{#1}}}

\author{Sören Gutzeit, Daniel Vogel, Artjom Siebert, Patrick Koteras}
\title{Tree-Height Balancing}

% Titlepage
\begin{document}
\begin{titlepage}
	\begin{center}
		\includegraphics[width=1.0\textwidth]{images/logo}\\
		\vspace{3cm}
		\begin{Huge}
			 Tree-Height Balacing
		\end{Huge} \linebreak
		Ausarbeitung im Modul Compilerbau 2 \\
		\vspace{3cm}
		Sören Gutzeit, Daniel Vogel, Artjom Siebert, Patrick Koteras \\
		\today
	\end{center}
\end{titlepage}

% Leere Seite nach Titlepage
\newpage 
\thispagestyle{empty}
\quad
\newpage

% Inhaltsverzeichnis
\pagenumbering{Roman}
\setcounter{page}{1}
\tableofcontents

% Sichern der aktuellen Seitenzahl des Vorspanns
\newcounter{last_roman}
\setcounter{last_roman}{\value{page}}
\newpage

%Seiten
\pagenumbering{arabic}
\chapter{Einführung}
\label{Einfuehrung}




\section{Nutzen und Erwartungen}
\label{Nutzen}
\section{Testing}
\label{Testing}

\subsection{Algorithmus}
\label{Algorithmus}

\subsection{Single- vs. Multi-Core}
\label{SMCore}

% Leere Seite nach Inhalt
\newpage 
\thispagestyle{empty}
\quad
\newpage

% Schluss
\pagenumbering{Roman}
\setcounter{page}{\value{last_roman}}
\chapter*{Abkürzungsverzeichnis}
\addcontentsline{toc}{chapter}{Abkürzungsverzeichnis}
% \ac{abk} für automatisch Abkürzung oder Beschreibung
% \acs{abk} für Abkürzung
% \acl{abk} für Beschreibung
% \acf{abk} für Abkürzung + Beschreibung
\begin{acronym}
\setlength{\itemsep}{-\parsep}
\acro{THB}{Tree-Height BalancingL}
\end{acronym}
\newpage
\listoffigures
\addcontentsline{toc}{chapter}{Abbildungsverzeichnis}
\newpage
\phantomsection 
\lstlistoflistings
\addcontentsline{toc}{chapter}{Listings}
\newpage
\phantomsection
\addcontentsline{toc}{chapter}{Literatur}
\nocite{*}
\printbibliography

\end{document}